\documentclass{article}
\usepackage{amsmath}
\usepackage{xcolor}
\begin{document}

\textbf{CAP 02 - ECUACIONES DIFERENCIALES DE ORDEN SUPERIOR}

\textbf{REPRESENTACION DE UN E.D. de Primer Orden}


\begin{gather*}
\frac{dy}{dx} = f(x,y) (A)  \rightarrow \text{ Forma General} \\
\frac{dy}{dx} = \frac{x^{2}+sen(x)}{e^{x}+y} \rightarrow     \\
\text{\colorbox[RGB]{248,231,28}{$M$}}\text{\colorbox[RGB]{248,231,28}{$($}}\text{\colorbox[RGB]{248,231,28}{$x,y$}}\text{\colorbox[RGB]{248,231,28}{$)$}}dx + \text{\colorbox[RGB]{126,211,33}{$N$}}\text{\colorbox[RGB]{126,211,33}{$($}}\text{\colorbox[RGB]{126,211,33}{$x,y$}}\text{\colorbox[RGB]{126,211,33}{$)$}}dy = 0 (B)  \rightarrow \text{ Forma Diferencial} \\
e^{x}dy = (x^{2}+sen(y))dx \\
(\text{\colorbox[RGB]{248,231,28}{$x$}}\text{\colorbox[RGB]{248,231,28}{$^{2}$}}\text{\colorbox[RGB]{248,231,28}{$+sen$}}\text{\colorbox[RGB]{248,231,28}{$($}}\text{\colorbox[RGB]{248,231,28}{$y$}}\text{\colorbox[RGB]{248,231,28}{$)$}})dx \text{\colorbox[RGB]{126,211,33}{$-e$}}\text{\colorbox[RGB]{126,211,33}{$^{x}$}}dy = 0
\end{gather*}




\textbf{VARIABLE SEPARABLE.- }Es posible separar las variables en un extremo\textbf{ }


\begin{gather*}
\frac{dy}{dx} = f(x,y)  \rightarrow  f(x,y)\text{ deberias ser una combinancion de factores que } \\
\text{se puede separar o se deberia poder buscar  alguna forma algebraica de formarlos}
\end{gather*}



\begin{gather*}
\vspace{0.5em} \\
M(x,y)dx + N(x,y)dy = 0  \rightarrow  A(x)dx + B(y)dy =0
\end{gather*}



\begin{gather*}
\frac{dy}{dx} =e^{-y}(2x-4)         y(5)=0 \\
\int e^{y} dy=\int (2x-4)dx +C \\
e^{y}  =x^{2}-4x +C  \\
y = ln(x^{2}-4x +C) \\
\text{para determinar la solucion especifica para y}(5)=0  \rightarrow  y(x) =f(...) \\
0 = ln(5^{2}-4*5 +C) \\
1  =25-20 +C  \\
C =-4 \\
y = ln(x^{2}-4x -4)
\end{gather*}






EED
\begin{gather*}
\frac{dy}{dx} = \frac{x^{2}+sen(x)}{e^{x}+y} \\
\vspace{0.5em} \\
y\frac{dy}{dx} =\frac{x^{2}+sen(x)}{e^{x}} \\
ydy =\frac{x^{2}+sen(x)}{e^{x}}\text{dx   } \rightarrow    _{:(} \\
\vspace{0.5em} \\
\vspace{0.5em} \\
\frac{dy}{dx} = \frac{x^{2}+sen(x)}{xy} =\frac{x^{2}}{xy}  +\frac{sen(x)}{xy}  \rightarrow  ok
\end{gather*}


\textbf{E.D. Homogenea }$\rightarrow$  se basa en un intercambio de variable para convertir la E.D. a una de variable separable, para ello se debe cumplir lo siguiente:


\begin{gather*}
\frac{dy}{dx} = f(x,y)\text{ si incluyes el operador '𝜆'} \\
\frac{dy}{dx} = f(𝜆x,𝜆y)  \rightarrow \text{ por operaciones algebraicas,} \\
\text{ deberias poder volver a la e.d. original} \\
\frac{dy}{dx} = f(x,y) 
\end{gather*}


ejemplo:


\begin{gather*}
\frac{dy}{dx} = \frac{x^{2}+y^{2}}{xy} = \frac{(𝜆x)^{2}+(𝜆y)^{2}}{𝜆x𝜆y} =\frac{𝜆^{2}x^{2}+𝜆^{2}y^{2}}{𝜆^{2}xy} = \frac{\text{\colorbox[RGB]{248,231,28}{$𝜆$}}\text{\colorbox[RGB]{248,231,28}{$^{2}$}}(x^{2}+y^{2})}{\text{\colorbox[RGB]{248,231,28}{$𝜆$}}\text{\colorbox[RGB]{248,231,28}{$^{2}$}}xy}
\end{gather*}


Cambio de Variable


\begin{gather*}
\frac{dy}{dx} = \frac{x^{2}+y^{2}}{xy} =\frac{x^{2}}{xy} +\frac{y^{2}}{xy} = \frac{x}{y}+\frac{y}{x} \\
\vspace{0.5em} \\
\text{por criterio asumo que }\frac{x}{y} =v \\
\frac{dy}{dx} = v +v^{-1} \\
x = vy    \rightarrow    (uv)' = u'v+uv' \\
\frac{dx}{dy} = \frac{dv}{dy}y+v \\
\vspace{0.5em} \\
----------------------- \\
y = xv \\
  \frac{dy}{dx} = v+x\frac{dv}{dx}   \rightarrow A \\
\vspace{0.5em} \\
x = yu \\
\frac{dx}{dy} = u+y\frac{du}{dy}    \rightarrow B
\end{gather*}



\begin{gather*}
\frac{dy}{dx} = \frac{x^{2}+y^{2}}{xy}\text{  aplicamos A} \\
\frac{dy}{dx} = \frac{x^{2}+(xv)^{2}}{xxv} = \frac{\text{\colorbox[RGB]{248,231,28}{$x$}}\text{\colorbox[RGB]{248,231,28}{$^{2}$}}+\text{\colorbox[RGB]{248,231,28}{$x$}}\text{\colorbox[RGB]{248,231,28}{$^{2}$}}v^{2}}{\text{\colorbox[RGB]{248,231,28}{$x$}}\text{\colorbox[RGB]{248,231,28}{$^{2}$}}v} \\
v+x\frac{dv}{dx}=\frac{1+v^{2}}{v} \\
x\frac{dv}{dx}=\frac{1+v^{2}}{v}-v = \frac{1+v^{2}-v^{2}}{v} =\frac{1}{v} \\
x\frac{dv}{dx}=\frac{1}{v} \\
\int vdv=\int \frac{1}{x}dx+C \\
ln(a)+ln(b) = ln(ab) \\
\frac{v^{2}}{2} = ln(x)+C  =   ln(x)+ln(C) = ln(Cx) \\
v = \sqrt{2ln(Cx)} \\
\vspace{0.5em} \\
 \rightarrow  K=2C \\
v^{2} = 2ln(x)+K \\
\frac{y^{2}}{x^{2}}= 2ln(x)+K \\
y^{2} =x^{2}2ln(x)+x^{2}K
\end{gather*}




ED Primer Orden Lineal


\begin{gather*}
\text{\colorbox[RGB]{248,231,28}{$\frac{dy}{dx}$}}\text{\colorbox[RGB]{248,231,28}{$ + P$}}\text{\colorbox[RGB]{248,231,28}{$($}}\text{\colorbox[RGB]{248,231,28}{$x$}}\text{\colorbox[RGB]{248,231,28}{$)$}}\text{\colorbox[RGB]{248,231,28}{$y = Q$}}\text{\colorbox[RGB]{248,231,28}{$($}}\text{\colorbox[RGB]{248,231,28}{$x$}}\text{\colorbox[RGB]{248,231,28}{$)$}}    \rightarrow  (A) \\
a_{1}(x)\frac{dy}{dx} + a_{0}(x)y = b(x) \\
\frac{a_{1}(x)}{a_{1}(x)}\frac{dy}{dx} + \frac{a_{0}(x)}{a_{1}(x)}y = \frac{b(x)}{a_{1}(x)} \\
P(x) =\frac{a_{0}(x)}{a_{1}(x)}     ; Q(x) = \frac{b(x)}{a_{1}(x)} \\
\text{Metodo de resolucion} \\
\frac{dy}{dx}u(x) + P(x)u(x)y = Q(x)u(x) = (yu)' \\
(vz)' = v'z+vz' \\
(yu)'=y'u+u'y \\
u'=Pu \\
\int \frac{u'}{u}=\int P \\
ln(u) = \int P(x) \\
u = e^{\int P(x)} \\
a)\text{ Calculamos un factor de integracion u}(x) \\
u(x) = e^{\int P(x)dx} \\
b)\text{ y calcular la solucion con:} \\
y = u(x) ^{-1}(\int Q(x)u(x) +c)
\end{gather*}


Ejercicio


\begin{gather*}
x\frac{dy}{dx}+3y= 4x^{2}-3x \\
\text{para llevar a }(A)\text{ , debemos dividir entre 'x'} \\
\frac{dy}{dx}+\frac{3}{x}y= 4x-3 \\
P(x) = \frac{3}{x};  Q(x)=4x-3 \\
\text{Calculamos el factor de integracion} \\
u(x) = e^{\int P(x)dx} \\
u(x) = e^{\int \frac{3}{x}dx}=e^{3\int \frac{1}{x}dx}= e^{3ln(x)} \\
a*ln(b)  = ln(b^{a}) \\
=e^{ln(x^{3})}=x^{3} \\
\text{determinamos la solucion y}(x) \\
y = u(x) ^{-1}(\int Q(x)u(x) +c) \\
y = x ^{-3}(\int (4x-3)x^{3} +c) \\
y = x ^{-3}(\int (4x^{4}-3x^{3} +c) \\
y = x ^{-3}( \frac{4}{5}x^{5}-\frac{3}{4}x^{4}+c) \\
y = \frac{4}{5}x^{2}-\frac{3}{4}x+x ^{-3}c
\end{gather*}


E.D. de Bernoulli


\begin{gather*}
\frac{dy}{dx} + P(x)y = Q(x)\text{\colorbox[RGB]{248,231,28}{$y$}}\text{\colorbox[RGB]{248,231,28}{$^{n}$}} \\
donde n = es real diferente 1, 0 \\
\text{el objetivo es tratar de convertir la E.D. en una lineal} \\
\vspace{0.5em} \\
\text{lo primero es hacer desaparecer y}^{n}\text{ del miembro de la derecha} \\
y^{-n}\frac{dy}{dx} + P(x)y^{-n}y = Q(x) \\
y^{-n}\frac{dy}{dx} + P(x)y^{1-n} = Q(x) \\
aplicamos cambio de variable z = y^{1-n} \\
y^{-n}\frac{dy}{dx} + P(x)z = Q(x) \\
derivamos z = y^{1-n} \\
\frac{dz}{dy} =(1-n)y^{-n}  \\
\frac{dz}{dy} =(1-n)y^{-n} \frac{dx}{dx} \\
\frac{dz}{dx} =(1-n)y^{-n} \frac{dy}{dx} \\
\text{Multplicar por}(1-n) \\
(1-n)y^{-n}\frac{dy}{dx} + (1-n)P(x)z = (1-n)Q(x) \\
\text{reemplazamos} \\
\frac{dz}{dx} + (1-n)P(x)z = (1-n)Q(x)
\end{gather*}


Ejercicio


\begin{gather*}
y'+\frac{4}{x}y = x^{3}y^{2} ;   y(2) = -1 \\
n = 2 \\
\text{Dividimos entre y}^{2} \\
y^{-2}y'+\frac{4}{x}y^{-1} = x^{3} \\
z =y^{-1} \\
\frac{dz}{dx} = -1y^{-2}\frac{dy}{dx} \\
\text{Multiplicamos la E.D. por }(-1) \\
- y^{-2}y'-\frac{4}{x}y^{-1} = -x^{3} \\
\text{hacemos los cambios respectivos} \\
\frac{dz}{dx}-\frac{4}{x}z =-x^{3} \\
\text{calculamos el factor de integracion u}(x) \\
u(x) = e^{\int p(x)dx} = e^{-4\int \frac{1}{x}dx}=x^{-4} \\
z(x) = u(x)^{-1}(\int Q(x)u(x)+C) \\
z(x) = x^{4}(\int -x^{3}x^{-4}+C)  \\
z(x) =x^{4}ln(x^{-1})+x^{4}C \\
y^{-1} = x^{4}ln(x^{-1})+x^{4}C \\
y = \frac{1}{x^{4}ln(x^{-1})+x^{4}C}
\end{gather*}


\textbf{E.D. Homogeneas}


\begin{gather*}
\frac{\partial M(x,y)}{\partial y} = \frac{\partial N(x,y)}{\partial x} \\
\text{Prueba de Homogeneidad} \\
\text{En caso de que se cumpla} \rightarrow \text{ podemos aplicar:} \\
\frac{\partial g(x,y)}{\partial x} = M(x,y)   \rightarrow  (A) \\
\frac{\partial g(x,y)}{\partial y} = N(x,y) \rightarrow  (B) \\
\text{donde g}(x,y)\text{ es la solucion de la E.D. } \rightarrow  g(x,y) = c \\
\text{Si despejamos }(A) \\
g(x,y) =\int  M(x,y)dx + h(y) \\
\text{Si despejamos }(B) \\
g(x,y) =\int  N(x,y)dy + v(x) \\
\text{Posteriormente se vuelve nuevamente a derivar } \\
\text{con respecto a la otra variable y se iguala a la funcion respectiva cruzada}
\end{gather*}


Ejercicios


\begin{gather*}
(2x+y)+(x+2y)y'=0; y(3)=1 \\
(2x+y)+(x+2y)\frac{dy}{dx}=0; \\
(2x+y)dx+(x+2y)dy=0 \\
\frac{\partial (2x+y)}{\partial y} = \frac{\partial (x+2y)}{\partial x} \\
1 =1  \\
\frac{\partial g(x,y)}{\partial x} =2x+y \\
g(x,y)= \int (2x+y)dx + h(y) \\
g(x,y)= x^{2}+xy+h(y) \\
\frac{\partial g(x,y)}{\partial y} = x+h'(y) = x+2y \\
h'(y) = 2y \\
h(y) = \int 2ydy = y^{2}+k \\
g(x,y)= x^{2}+xy+y^{2}+k = c \\
g(x,y)= x^{2}+xy+y^{2}+A ; donde A = c-k \\
\vspace{0.5em} \\
\text{utilizando las formulas al reves} \\
\frac{\partial g(x,y)}{\partial y} = x+2y \\
g(x,y) = \int (x+2y)dy + v(x) \\
g(x,y) = xy+y^{2}+ v(x) \\
\frac{\partial g(x,y)}{\partial x} =y+ v'(x) = M(x,y)=2x+y \\
v(x) = x^{2}+c \\
g(x,y) = xy+y^{2}+ x^{2} =A; donde A = c-k \\
\text{calculamos la solucion particular con } \\
y(x=3)=1   (x,y)  \rightarrow  (3,1) \\
(3)(1)+(1)^{2}+ (3)^{2} =A \\
A = 13 \\
g(x,y) = xy+y^{2}+ x^{2} =13
\end{gather*}







\begin{gather*}
ycosx+ye^{xy}+(senx+xe^{xy})y' = 0
\end{gather*}





\begin{gather*}
y^{2}+x^{3}+xyy'=0
\end{gather*}





\end{document}
